% GENERATED FILE. Edit canonical lessons, then run npm run generate:books.
% canonical-source-hash: fnv1a64:c621fc113734e46e
% canonical-lessons: GU-C496-dhavvu, GU-C496-hicolvu, GU-C496-namkaran, GU-C496-vadhavvu, GU-C496-balotiyu

\chapter{A New Baby}
\label{ch:gu-a-new-baby}
\hlchaptermodality{496}

\begin{chapteropening}
\textbf{By the end of this chapter:} \emph{I can talk about the first days of a new baby.}

Complete the last lesson of chapter 496: I can talk about the first days of a new baby.
\end{chapteropening}

\section[\texorpdfstring{\gu{ધાવવું} (dhāvvũ)}{ધાવવું (dhāvvũ)}]{\gu{ધાવવું} (dhāvvũ) — to suckle; to feed at the breast}

\label{lesson:GU-C496-dhavvu}



\begin{quote}
Before the new one: say the Gujarati for to come into the world, then the Gujarati for to give birth to.
\end{quote}

\subsection*{You'll want to know: \gu{ધાવવું}}
\textbf{\gu{ધાવવું}} — \emph{dhāvvũ} — \textquotedblleft{}to suckle; to feed at the breast\textquotedblright{}.

\begin{quote}
\gu{બાળક} \gu{માતાને} \gu{ધાવે} \gu{છે}.
\end{quote}

\begin{quote}
\emph{bāḷak mātāne dhāve che.}
\end{quote}

\begin{quote}
The baby is feeding at its mother's breast.
\end{quote}

Past (it …): \gu{તે} \gu{બાળક} \gu{ધાવ્યું}, \emph{te bāḷak dhāvyũ}.

\begin{cousinweb}
\textbf{\gu{ધાવવું}} (\emph{dhāvvũ}) is inherited from Sanskrit \emph{dhāpayatē}, \textquotedblleft{}suckles\textquotedblright{}; CDIAL lists the Gujarati verb, \textquotedblleft{}to suck at the breast\textquotedblright{} (CDIAL 6784).
\end{cousinweb}

\begin{grammarlens}[title={What you've built}]
One more word for describing what you lived and felt.
\end{grammarlens}

\subsection*{Your turn}
\begin{itemize}
\raggedright
  \item \emph{Say it:} \emph{dhāvvũ}
  \item \emph{Say it:} \emph{dhāvvũ}, once more
  \item \emph{Say it:} \emph{bāḷak mātāne dhāve che.} — The baby is feeding at its mother's breast.
  \item \emph{Recall:} say the Gujarati for to come into the world, then the Gujarati for to give birth to, then say \emph{dhāvvũ} again
\end{itemize}

\begin{tcolorbox}[breakable,colback=teal!4,colframe=teal!35!black,title={Before you move on}]
What does \gu{ધાવવું} mean? (\textquotedblleft{}To suckle; to feed at the breast\textquotedblright{}.) Say it once more.
\end{tcolorbox}
\section[\texorpdfstring{\gu{હીંચોળવું} (hī̃coḷvũ)}{હીંચોળવું (hī̃coḷvũ)}]{\gu{હીંચોળવું} (hī̃coḷvũ) — to rock a cradle; to rock a swing}

\label{lesson:GU-C496-hicolvu}



\begin{quote}
Before the new one: say the Gujarati for to give birth to, then the Gujarati for to suckle.
\end{quote}

\subsection*{You'll want to know: \gu{હીંચોળવું}}
\textbf{\gu{હીંચોળવું}} — \emph{hī̃coḷvũ} — \textquotedblleft{}to rock a cradle; to rock a swing\textquotedblright{}.

\begin{quote}
\gu{દાદી} \gu{ગીત} \gu{ગાઈને} \gu{બાળકને} \gu{હિંચકામાં} \gu{હીંચોળે} \gu{છે}.
\end{quote}

\begin{quote}
\emph{dādī gīt gāīne bāḷakne hinckāmā̃ hī̃coḷe che.}
\end{quote}

\begin{quote}
Grandmother sings a song and rocks the baby on the swing.
\end{quote}

Past, with \gu{તેણે} (he or she …): \gu{તેણે} \gu{બાળકને} \gu{હીંચોળ્યું}, \emph{teṇe bāḷakne hī̃coḷyũ}.

\begin{cousinweb}
\textbf{\gu{હીંચોળવું}} (\emph{hī̃coḷvũ}) is built on the stem that CDIAL lists as Gujarati \emph{hĩcvũ}, \textquotedblleft{}to swing\textquotedblright{}, under a reconstructed \emph{*hicc-} (CDIAL 14083); \textbf{\gu{હિંચકો}} (\emph{hincko}), \textquotedblleft{}a swing\textquotedblright{}, is its cousin.
\end{cousinweb}

\begin{grammarlens}[title={What you've built}]
One more word for describing what you lived and felt.
\end{grammarlens}

\subsection*{Your turn}
\begin{itemize}
\raggedright
  \item \emph{Say it:} \emph{hī̃coḷvũ}
  \item \emph{Say it:} \emph{hī̃coḷvũ}, once more
  \item \emph{Say it:} \emph{dādī gīt gāīne bāḷakne hinckāmā̃ hī̃coḷe che.} — Grandmother sings a song and rocks the baby on the swing.
  \item \emph{Recall:} say the Gujarati for to give birth to, then the Gujarati for to suckle, then say \emph{hī̃coḷvũ} again
\end{itemize}

\begin{tcolorbox}[breakable,colback=teal!4,colframe=teal!35!black,title={Before you move on}]
What does \gu{હીંચોળવું} mean? (\textquotedblleft{}To rock a cradle; to rock a swing\textquotedblright{}.) Say it once more.
\end{tcolorbox}
\section[\texorpdfstring{\gu{નામકરણ} (nāmkaraṇ)}{નામકરણ (nāmkaraṇ)}]{\gu{નામકરણ} (nāmkaraṇ) — a naming ceremony}

\label{lesson:GU-C496-namkaran}



\begin{quote}
Before the new one: say the Gujarati for to suckle, then the Gujarati for to rock a cradle.
\end{quote}

\subsection*{You'll want to know: \gu{નામકરણ}}
\textbf{\gu{નામકરણ}} — \emph{nāmkaraṇ} — \textquotedblleft{}a naming ceremony\textquotedblright{}.

\begin{quote}
\gu{નામકરણમાં} \gu{ફોઈએ} \gu{બાળકનું} \gu{નામ} \gu{નક્કી} \gu{કર્યું}.
\end{quote}

\begin{quote}
\emph{nāmkaraṇmā̃ phoīe bāḷaknũ nām nakkī karyũ.}
\end{quote}

\begin{quote}
At the naming ceremony the baby's paternal aunt chose its name.
\end{quote}

A neuter noun.

\begin{cousinweb}
\textbf{\gu{નામકરણ}} (\emph{nāmkaraṇ}) is a tatsama, Sanskrit \emph{nāmakaraṇa}, \textquotedblleft{}the giving of a name\textquotedblright{}, as Monier-Williams records it.
\end{cousinweb}

\begin{grammarlens}[title={What you've built}]
One more word for describing what you lived and felt.
\end{grammarlens}

\subsection*{Your turn}
\begin{itemize}
\raggedright
  \item \emph{Say it:} \emph{nāmkaraṇ}
  \item \emph{Say it:} \emph{nāmkaraṇ}, once more
  \item \emph{Say it:} \emph{nāmkaraṇmā̃ phoīe bāḷaknũ nām nakkī karyũ.} — At the naming ceremony the baby's paternal aunt chose its name.
  \item \emph{Recall:} say the Gujarati for to suckle, then the Gujarati for to rock a cradle, then say \emph{nāmkaraṇ} again
\end{itemize}

\begin{tcolorbox}[breakable,colback=teal!4,colframe=teal!35!black,title={Before you move on}]
What does \gu{નામકરણ} mean? (\textquotedblleft{}A naming ceremony\textquotedblright{}.) Say it once more.
\end{tcolorbox}
\section[\texorpdfstring{\gu{વધાવવું} (vadhāvvũ)}{વધાવવું (vadhāvvũ)}]{\gu{વધાવવું} (vadhāvvũ) — to welcome with rites; to greet with joy}

\label{lesson:GU-C496-vadhavvu}



\begin{quote}
Before the new one: say the Gujarati for to rock a cradle, then the Gujarati for a naming ceremony.
\end{quote}

\subsection*{You'll want to know: \gu{વધાવવું}}
\textbf{\gu{વધાવવું}} — \emph{vadhāvvũ} — \textquotedblleft{}to welcome with rites; to greet with joy\textquotedblright{}.

\begin{quote}
\gu{દાદીએ} \gu{બાળકને} \gu{ગીત} \gu{ગાઈને} \gu{વધાવ્યું}.
\end{quote}

\begin{quote}
\emph{dādīe bāḷakne gīt gāīne vadhāvyũ.}
\end{quote}

\begin{quote}
Grandmother welcomed the baby with a song.
\end{quote}

Past, with \gu{તેણે} (he or she …): \gu{તેણે} \gu{બાળકને} \gu{વધાવ્યું}, \emph{teṇe bāḷakne vadhāvyũ}.

\begin{cousinweb}
\textbf{\gu{વધાવવું}} (\emph{vadhāvvũ}) is inherited from Sanskrit \emph{vardhāpayati}, \textquotedblleft{}increases, gladdens, congratulates\textquotedblright{}; CDIAL lists the Gujarati verb, \textquotedblleft{}to welcome ceremonially\textquotedblright{} (CDIAL 11383).
\end{cousinweb}

\begin{grammarlens}[title={What you've built}]
One more word for describing what you lived and felt.
\end{grammarlens}

\subsection*{Your turn}
\begin{itemize}
\raggedright
  \item \emph{Say it:} \emph{vadhāvvũ}
  \item \emph{Say it:} \emph{vadhāvvũ}, once more
  \item \emph{Say it:} \emph{dādīe bāḷakne gīt gāīne vadhāvyũ.} — Grandmother welcomed the baby with a song.
  \item \emph{Recall:} say the Gujarati for to rock a cradle, then the Gujarati for a naming ceremony, then say \emph{vadhāvvũ} again
\end{itemize}

\begin{tcolorbox}[breakable,colback=teal!4,colframe=teal!35!black,title={Before you move on}]
What does \gu{વધાવવું} mean? (\textquotedblleft{}To welcome with rites; to greet with joy\textquotedblright{}.) Say it once more.
\end{tcolorbox}
\section[\texorpdfstring{\gu{બાળોતિયું} (bāḷotiyũ)}{બાળોતિયું (bāḷotiyũ)}]{\gu{બાળોતિયું} (bāḷotiyũ) — a baby's nappy; a cloth for a baby}

\label{lesson:GU-C496-balotiyu}



\begin{quote}
Before the new one: say the Gujarati for a naming ceremony, then the Gujarati for to welcome with rites.
\end{quote}

\subsection*{You'll want to know: \gu{બાળોતિયું}}
\textbf{\gu{બાળોતિયું}} — \emph{bāḷotiyũ} — \textquotedblleft{}a baby's nappy; a cloth for a baby\textquotedblright{}.

\begin{quote}
\gu{માતાએ} \gu{બાળકનું} \gu{ભીનું} \gu{બાળોતિયું} \gu{બદલ્યું}.
\end{quote}

\begin{quote}
\emph{mātāe bāḷaknũ bhīnũ bāḷotiyũ badalyũ.}
\end{quote}

\begin{quote}
The mother changed the baby's wet nappy.
\end{quote}

A neuter noun.

\begin{cousinweb}
\textbf{\gu{બાળોતિયું}} (\emph{bāḷotiyũ}) goes back to a reconstructed \emph{*bāladhōtta-}, \textquotedblleft{}child's cloth\textquotedblright{}, built from Sanskrit \emph{bāla-}, \textquotedblleft{}child\textquotedblright{}, and \emph{*dhōtta-}, \textquotedblleft{}cloth\textquotedblright{}; CDIAL lists the Gujarati word, \textquotedblleft{}cloth for a baby\textquotedblright{} (CDIAL 9219).
\end{cousinweb}

\begin{grammarlens}[title={What you've built}]
One more word for describing what you lived and felt.
\end{grammarlens}

\subsection*{Your turn}
\begin{itemize}
\raggedright
  \item \emph{Say it:} \emph{bāḷotiyũ}
  \item \emph{Say it:} \emph{bāḷotiyũ}, once more
  \item \emph{Say it:} \emph{mātāe bāḷaknũ bhīnũ bāḷotiyũ badalyũ.} — The mother changed the baby's wet nappy.
  \item \emph{Recall:} say the Gujarati for a naming ceremony, then the Gujarati for to welcome with rites, then say \emph{bāḷotiyũ} again
\end{itemize}

\begin{tcolorbox}[breakable,colback=teal!4,colframe=teal!35!black,title={Before you move on}]
What does \gu{બાળોતિયું} mean? (\textquotedblleft{}A baby's nappy; a cloth for a baby\textquotedblright{}.) Say it once more.
\end{tcolorbox}
